% Part: normal-modal-logic
% Chapter: syntax-and-semantics
% Section: language-modal-logic

\documentclass[../../../include/open-logic-section]{subfiles}

\begin{document}

\olfileid[fr]{nml}{syn}{lan}

\olsection{Le langage de la logique modale de base}

\begin{defn}
Le langage de base de la logique modale comprend :
\begin{enumerate}
  \tagitem{prvFalse}{La constante propositionnelle pour la !!{falsity}~$\lfalse$.}{}
  \tagitem{prvTrue}{La constante propositionnelle pour la !!{truth}~$\ltrue$.}{}
  \item Un ensemble !!{denumerable} de !!{propositional variable}s : $\Obj
    p_0$, $\Obj p_1$, $\Obj p_2$, \dots
  \item Les connecteurs propositionnels : \startycommalist
  \iftag{prvNot}{\ycomma $\lnot$ (négation)}{}%
  \iftag{prvAnd}{\ycomma $\land$ (conjonction)}{}%
  \iftag{prvOr}{\ycomma $\lor$ (disjonction)}{}%
  \iftag{prvIf}{\ycomma $\lif$ (!!{conditional})}{}%
  \iftag{prvIff}{\ycomma $\liff$ (!!{biconditional})}{}.
  \tagitem{prvBox}{L'opérateur modal $\Box$.}{}
  \tagitem{prvDiamond}{L'opérateur modal $\Diamond$.}{}
\end{enumerate}
\end{defn}

\begin{defn}
Les \emph{!!^{formula}s} du langage modal de base sont définies
  inductivement comme suit :
\begin{enumerate}
\tagitem{prvFalse}{$\lfalse$ est une !!{formula} atomique.}{}

\tagitem{prvTrue}{$\ltrue$ est une !!{formula} atomique.}{}

\item Toute !!{propositional variable} $\Obj p_i$ est une !!{formula}
  atomique.

\tagitem{prvNot}{Si $!A$ est !!a{formula}, alors $\lnot !A$ est
  !!a{formula}.}{}

\tagitem{prvAnd}{Si $!A$ et $!B$ sont des !!{formula}s, alors
  $(!A \land !B)$ est !!a{formula}.}{}

\tagitem{prvOr}{Si $!A$ et $!B$ sont des !!{formula}s, alors
  $(!A \lor !B)$ est !!a{formula}.}{}

\tagitem{prvIf}{Si $!A$ et $!B$ sont des !!{formula}s, alors
  $(!A \lif !B)$ est !!a{formula}.}{}

\tagitem{prvIff}{Si $!A$ et $!B$ sont des !!{formula}s, alors
  $(!A \liff !B)$ est !!a{formula}.}{}

\tagitem{prvBox}{Si $!A$ est !!a{formula}, alors $\Box !A$ est
  !!a{formula}.}{}

\tagitem{prvDiamond}{Si $!A$ est !!a{formula}, alors $\Diamond !A$
  est !!a{formula}.}{}

\tagitem{limitClause}{Rien d'autre n'est !!a{formula}.}{}
\end{enumerate}
\end{defn}

\begin{tagblock}{defNot,defOr,defAnd,defIf,defIff,defTrue,defFalse,defBox,defDiamond}
\begin{defn}
Les formules construites au moyen des opérateurs définis s'entendent
selon les abréviations suivantes :

\begin{tagenumerate}{defTrue,defFalse,defNot,defOr,defAnd,defIf,defIff,defBox,defDiamond}

\tagitem{defTrue}{$\ltrue$ abrège
  \iftag{prvFalse}{$\lnot\lfalse$}{$(!A \lor \lnot !A)$ pour une
    !!{formula} atomique~$!A$ fixée}.}{}

\tagitem{defFalse}{$\lfalse$ abrège
  \iftag{prvTrue}{$\lnot\ltrue$}{$(!A \land \lnot !A)$ pour une
    !!{formula} atomique~$!A$ fixée}.}{}

\tagitem{defNot}{$\lnot !A$ abrège $!A \lif \lfalse$.}{}

\tagitem{defOr}{$!A \lor !B$ abrège
  \iftag{prvAnd}{$\lnot(\lnot !A \land \lnot !B)$}{$\lnot !A \lif
    !B$}.}{}

\tagitem{defAnd}{$!A \land !B$ abrège
  \iftag{prvOr}{$\lnot(\lnot !A \lor \lnot !B)$}{$\lnot (!A \lif
    \lnot !B)$}.}{}

\tagitem{defIf}{$!A \lif !B$ abrège
  \iftag{prvOr}{$(\lnot !A \lor !B)$}{$\lnot (!A \land \lnot !B)$}.}{}

\tagitem{defIff}{$!A \liff !B$ abrège $(!A \lif !B) \land (!B
  \lif !A)$.}{}

\tagitem{defBox}{$\Box !A$ abrège $\lnot\Diamond\lnot !A$. }{}

\tagitem{defDiamond}{$\Diamond !A$ abrège $\lnot\Box\lnot !A$.}{}
\end{tagenumerate}
\end{defn}
\end{tagblock}

Si !!a{formula}~$!A$ ne contient ni $\Box$ ni $\Diamond$, nous
disons qu'elle est \emph{sans opérateur modal}.

\end{document}
